\documentclass[12pt]{article}
\usepackage{amsmath,amssymb}

\begin{document}
\section*{{2020 AIME II Problems/Problem 5}}
Source: AoPS Wiki

\subsection*{Solution}
Let's work backwards. The minimum base-sixteen representation of $g(n)$ that cannot be expressed using only the digits $0$ through $9$ is $A_{16}$ , which is equal to $10$ in base 10. Thus, the sum of the digits of the base-eight representation of the sum of the digits of $f(n)$ is $10$ . The minimum value for which this is achieved is $37_8$ . We have that $37_8 = 31$ . Thus, the sum of the digits of the base-four representation of $n$ is $31$ . The minimum value for which this is achieved is $13,333,333,333_4$ . We just need this value in base 10 modulo 1000. We get $13,333,333,333_4 = 3(1 + 4 + 4^2 + \dots + 4^8 + 4^9) + 4^{10} = 3\left(\dfrac{4^{10} - 1}{3}\right) + 4^{10} = 2*4^{10} - 1$ . Taking this value modulo $1000$ , we get the final answer of $\boxed{151}$ . (If you are having trouble with this step, note that $2^{10} = 1024 \equiv 24 \pmod{1000}$ ) ~ TopNotchMath

\end{document}
